\documentclass{zhvt-classic}

% 使用本地源流明体（GenRyuMin2 TW）按文件路径加载，不依赖系统字体名
% 版框（内框／界栏／书口侧栏／鱼尾）改用与版心印章同一印泥棕色：
%   印章图片 yshs.png 实色 = RGB(135,68,52) = HTML 874434
\usepackage{xcolor}
\usepackage{tikz}
% sealbrown 双注册：xcolor 供 TikZ/pgf，L3 供 cls 版框
\definecolor{sealbrown}{HTML}{874434}
\ExplSyntaxOn
\color_set:nnn { sealbrown } { HTML } { 874434 }
\ExplSyntaxOff

\zhvtset{
  grid color       = sealbrown,
  main font        = GenRyuMin2TW-SB,
  main font path   = C:/Users/Administrator/texmf/fonts/opentype/public/GenRyuMin2TW/,
jiazhu font      = GenRyuMin2TW-M,
  jiazhu font path = C:/Users/Administrator/texmf/fonts/opentype/public/GenRyuMin2TW/,
  % 页码框容纳 4 字：正常页「九十七」+前置空 1 格=4 行恰好；
  % 100+ 页「一百零一」共 5 行超出，末字「一/二/三」被框裁掉（2026-10-04 实测），故放宽到 5
  shukou page number chars = 5,
  % 书口书名框高从默认 6 字减到 4 字，把省出的 2 字高拨给 chapter 框，
  % 使 6 字书名+卷名放进 9 字高的 chapter 框后不再压线（2026-10-05）
  shukou title chars = 4,
}

\title{史记}[孔子世家]
\author{汉·司马迁}
\maker{闲叟漫记·重製}


% ===== 索隐述赞样式（2026-11-08 四次定稿 · 愚大人亲手改定参数）=====
% ⚠ 竖排前提：自定义字体家族必须带 RawFeature={vertical:+vert}；
%   simhei.ttf 加 +vhal 会整体加载失败静默回退，只写 +vert；simfang 同理。
\newCJKfontfamily{\shuzanfont}[Path=C:/Windows/Fonts/,RawFeature={vertical:+vert}]{simhei.ttf}
\newCJKfontfamily{\shuzanbody}[Path=C:/Windows/Fonts/,RawFeature={vertical:+vert}]{simfang.ttf}
% 题头：152x30pt 圆角 6pt 棕底白字；baseline (0,2) 使整框落入界行格子不压版框线；
% 白字节点 x 不补偿、y +8.2pt（像素级墨迹实测 dx/dy≈0，L/R 留白约 20pt）。
\newcommand{\shuzanhead}[1]{%
  \begin{tikzpicture}[baseline={(0,2)}]
    \fill[sealbrown, rounded corners=6pt]
      (-76pt,-15pt) rectangle (76pt,15pt);
    \node[text=white, inner sep=0pt, anchor=center] at (0pt,8.2pt)
      {\shuzanfont\fontsize{23}{23}\selectfont #1};
  \end{tikzpicture}%
}
% ★ 修正书口卷名位置：chapter 框起始位置改回 cls 默认值，避免全局影响。
%   只在每个卷首页通过 \juan 命令里的空命令微调该页位置（安全，不会跨页污染）。
\ExplSyntaxOn
\AtBeginDocument{
  \dim_set:Nn \l__zhvt_shukou_chapter_title_start_dim
    { \l__zhvt_shukou_title_height_dim + \l__zhvt_font_size_dim * 2 }
  \dim_set:Nn \l__zhvt_shukou_page_number_start_dim
    { \l__zhvt_shukou_chapter_title_start_dim + \l__zhvt_shukou_chapter_title_height_dim }
}
\ExplSyntaxOff

\begin{document}

% ★★ 版心（书口）排布覆写：把书名放到鱼尾之下（2026-10-04）
%
%   cls 默认版心自上而下：书名（title 静态框，在鱼尾之上）→ 鱼尾 → 篇名（chapter 动态框）→ 页码
%   本模板改为：          鱼尾（位置不动）→ 书名 → 篇名 → 页码分割线 → 页码
%
%   手法：把 cls 放在 title 静态框里的书名，改放进 chapter 动态框（鱼尾之下）；
%         title 静态框留空。书名取 \l__zhvt_title_str，即 \title{...} 的第一参数，
%         所以换书时只改 \title，不必在此重复写书名。
%
%   ⚠ 三条必须遵守的写法（2026-10-04 实测）：
%     ① 汉字必须用 \hbox{...} 包住。flowfram 会把框内容装进一个宽 14.4mm（= 一个 line_sep
%        = 书口栏宽）的纵向 minipage 并按该宽度折行，而整页靠 PDF /Rotate 90 显示、
%        框内文字本身不变向 —— 不包 \hbox 时「围炉夜话」4 字会被逐字折成 4 行，
%        在版面上就成了横排 4 字并横向溢出版框。
%     ② 竖排下**一个 \hbox 就是一整列**。书名与篇名要同处一列，必须放进同一个 \hbox；
%        分成两个 \hbox 时，后一个会被当成新的一列、甩到隔壁栏去
%        （实测「序」跑到书口栏外侧 14.4mm 的正文列上）。中间的 \quad 在框内是横向空白，
%        显示出来正好是书名与篇名之间的纵向间隔。
%     ③ ★ 内容盒要写成 **与框等宽的定宽盒** `\hbox to \l__zhvt_shukou_chapter_title_height_dim{…\hfil}`。
%        实测：chapter 框是把内容盒沿纵向**居中**摆放的，内容宽度一变，书名就跟着漂：
%        自然宽度的 \hbox（书名 38.4mm）落在 x=157.2mm，加篇名后（57.6mm）落在 x=152.4mm，
%        目录页（无篇名）与正文页的书名就差半格。做成等宽盒后，无论有没有篇名，
%        内容固定从框顶起排，实测两页均落在 x=142.8mm，不再漂。
\ExplSyntaxOn
% 把书口 chapter 框（书名+篇名）起始位置下移 1 字高，避免 4 字卷名顶到鱼尾
% ⚠ cls 的 \__zhvt_calculate_dimensions: 在 \AtBeginDocument 里执行，会重算该 dim；
%   必须把覆写挂到 begin document 之后，否则会被二次覆盖（2026-10-05 实测）。
% (旧方案已移到导言区，直接修改 chapter 框起始位置)
% 卷标题命令 #1=书口篇名(鱼尾) #2=正文卷名（含卷序，卷名内空格用 ~）
%
% ★★ 目录页两项修订（2026-10-05，愚大人指出）★★
%   ① 页码起点统一对齐：
%      原状 —— 目录项写成 \par\hspace{4em}卷名\hspace{4em}页码，页码起点 = 卷名实宽 + 4em。
%      卷名实宽随字数变：「卷一~梁惠王上」6 字 +~、「卷七~离娄上」5 字 +~，
%      差 1 个字位 → 卷七~卷十 4 条的页码整体比其余 10 条高 27.21pt（一个字），看着上下不齐。
%      改法 —— 卷名补全角空格（宽 = 1 字 = 27.215pt）补足到 **7 字符位**再写进目录，
%      14 条页码起点遂全部落在同一条竖线上（起点对齐；末字随页码字数自然参差）。
%      ⚠ 只对目录里的副本 \l_tmpc_tl 补，正文卷名仍用 \l_tmpa_tl 原样 —— 正文像素不受影响。
%   ② 卷名整体上移半字位：前导 \hspace 由 4em 改 3.5em。
%      实测原状目录块在版框内偏下（天头留白 139.6pt / 地脚留白 110.7pt），
%      上移 0.5em = 13.6pt 后上下基本均分，视觉居中。
\NewDocumentCommand { \juan } { mm } {
	\clearpage
	\tl_set:Nn \l__zhvt_chapter_title_tl { #1 }
	\tl_set:Nn \l_tmpa_tl { #2 }
	\tl_set_eq:NN \l_tmpc_tl \l_tmpa_tl
	\int_step_inline:nn { 7 - \str_count:n { #2 } }
	  { \tl_put_right:Nn \l_tmpc_tl { 　 } }
	\tl_set:Nx \l_tmpb_tl { \zhnum { page } }
	\tl_put_left:Nn \l_tmpb_tl { \hspace { 4em } }
	\tl_put_left:Nn \l_tmpb_tl { \l_tmpc_tl }
	\tl_put_left:Nn \l_tmpb_tl { \par \hspace { 3.5em } }
	\addtocontents { toc } { \l_tmpb_tl }
	\skip_horizontal:n { \l__zhvt_font_size_dim * 4 }
	\l_tmpa_tl
	\par
	\thispagestyle { empty }
}
\cs_set_protected:Npn \__zhvt_set_shukou_content: {
	% ★ 鱼尾上方留白（title 静态框）嵌入闲叟钤印：
	%   整页 /Rotate 90，故此处把图用 \rotatebox{90} 转正为竖章；
	%   尺寸卡在版心条幅内（书内约 7mm 宽 × 33.8mm 高，留余量不压鱼尾）。
	\setstaticcontents * { title } {
		\vbox_to_ht:nn { \l__zhvt_baseline_skip_dim } {
			\vskip 14.2pt
			\hbox_to_wd:nn { \l__zhvt_shukou_title_height_dim } {
				\hfil
				\rotatebox [ origin = c ] { 90 } { \includegraphics [ width = 7mm , height = 33.8mm ] { yshs.png } }
				\hfil
			}
		}
	}
% ★ 版心（书口）统一改法（2026-11-16 愚大人定）：18pt＋框内居中＋书口文字同版框色
%   ① 原 27.2pt 时 7 字长篇名（如「廉颇蔺相如列传」）末字溢出 chapter 框、压住页码；
%      改 \fontsize{18pt}{18pt}，前后 \hfil 令字串在框内居中（不随篇名字数漂移）。
%   ② \vbox{\vskip 8.2pt} 定量下移，抵消缩字号后栏宽方向 8.2pt 上偏（像素实测 y 中心 543.3 = 原版）。
%      ⚠ \raisebox / \vspace* 在 flowfram 框内无效；换字号须同步换补偿：16pt→10.1、19pt→7.4。
%   ③ chapter 与 page~number 均置 \color{sealbrown}（HTML 874434），与版框／界行同色；正文仍黑。
%   ⚠ 页码字位、钤印、鱼尾、界行、正文一律未动；五部合订本同此口径。
	\setdynamiccontents * { chapter }
	{ \color{sealbrown} \vbox { \vskip 8.2pt \hbox to \l__zhvt_shukou_chapter_title_height_dim
		{ \hfil \fontsize{18pt}{18pt}\selectfont \l__zhvt_title_str \quad \l__zhvt_chapter_title_tl \hfil } } }
	\setdynamiccontents * { page~number } { \color{sealbrown} \quad \zhnum { page } }
}
\__zhvt_set_shukou_content:
\ExplSyntaxOff

\maketitle

% 单篇，无多卷目录，不渲染目录页

\juan{孔子世家}{卷四十七~孔子世家第十七}
\hspace*{2em}孔子生鲁昌平乡陬邑。其先宋人也。曰孔防叔。防叔生伯夏。伯夏生叔梁纥。纥与颜氏女野合而生孔子。祷於尼丘得孔子。鲁襄公二十二年而孔子生。生而首上圩顶。故因名曰丘云。字仲尼。姓孔氏。\\
\hspace*{2em}丘生而叔梁纥死。葬於防山。防山在鲁东。由是孔子疑其父墓处。母讳之也。孔子为兒嬉戏。常陈俎豆。设礼容。孔子母死。乃殡五父之衢。盖其慎也。陬人輓父之母诲孔子父墓。然後往合葬於防焉。\\
\hspace*{2em}孔子要绖。季氏飨士。孔子与往。阳虎绌曰。﹁季氏飨士。非敢飨子也。﹂孔子由是退。\\
\hspace*{2em}孔子年十七。鲁大夫孟釐子病且死。诫其嗣懿子曰。﹁孔丘。圣人之後。灭於宋。其祖弗父何始有宋而嗣让厉公。及正考父佐戴。武。宣公。三命兹益恭。故鼎铭云。﹁一命而偻。再命而伛。三命而俯。循墙而走。亦莫敢余侮。饘於是。粥於是。以餬余口。﹂其恭如是。吾闻圣人之後。虽不当世。必有达者。今孔丘年少好礼。其达者欤。吾即没。若必师之。﹂及釐子卒。懿子与鲁人南宫敬叔往学礼焉。是岁。季武子卒。平子代立。\\
\hspace*{2em}孔子贫且贱。及长。尝为季氏史。料量平。尝为司职吏而畜蕃息。由是为司空。已而去鲁。斥乎齐。逐乎宋。卫。困於陈蔡之间。於是反鲁。孔子长九尺有六寸。人皆谓之﹁长人﹂而异之。鲁复善待。由是反鲁。\\
\hspace*{2em}鲁南宫敬叔言鲁君曰。﹁请与孔子適周。﹂鲁君与之一乘车。两马。一竖子俱。適周问礼。盖见老子云。辞去。而老子送之曰。﹁吾闻富贵者送人以财。仁人者送人以言。吾不能富贵。窃仁人之号。送子以言。曰。﹁聪明深察而近於死者。好议人者也。博辩广大危其身者。发人之恶者也。为人子者毋以有己。为人臣者毋以有己。﹂﹂孔子自周反于鲁。弟子稍益进焉。\\
\hspace*{2em}是时也。晋平公淫。六卿擅权。东伐诸侯。楚灵王兵彊。陵轹中国。齐大而近於鲁。鲁小弱。附於楚则晋怒。附於晋则楚来伐。不备於齐。齐师侵鲁。\\
\hspace*{2em}鲁昭公之二十年。而孔子盖年三十矣。齐景公与晏婴来適鲁。景公问孔子曰。﹁昔秦穆公国小处辟。其霸何也。﹂对曰。﹁秦。国虽小。其志大。处虽辟。行中正。身举五羖。爵之大夫。起累绁之中。与语三日。授之以政。以此取之。虽王可也。其霸小矣。﹂景公说。\\
\hspace*{2em}孔子年三十五。而季平子与郈昭伯以斗鸡故得罪鲁昭公。昭公率师击平子。平子与孟氏。叔孙氏三家共攻昭公。昭公师败。奔於齐。齐处昭公乾侯。其後顷之。鲁乱。孔子適齐。为高昭子家臣。欲以通乎景公。与齐太师语乐。闻韶音。学之。三月不知肉味。齐人称之。\\
\hspace*{2em}景公问政孔子。孔子曰。﹁君君。臣臣。父父。子子。﹂景公曰。﹁善哉。信如君不君。臣不臣。父不父。子不子。虽有粟。吾岂得而食诸。﹂他日又复问政於孔子。孔子曰。﹁政在节财。﹂景公说。将欲以尼谿田封孔子。晏婴进曰。﹁夫儒者滑稽而不可轨法。倨傲自顺。不可以为下。崇丧遂哀。破产厚葬。不可以为俗。游说乞贷。不可以为国。自大贤之息。周室既衰。礼乐缺有间。今孔子盛容饰。繁登降之礼。趋详之节。累世不能殚其学。当年不能究其礼。君欲用之以移齐俗。非所以先细民也。﹂後景公敬见孔子。不问其礼。异日。景公止孔子曰。﹁奉子以季氏。吾不能。﹂以季孟之间待之。齐大夫欲害孔子。孔子闻之。景公曰。﹁吾老矣。弗能用也。﹂孔子遂行。反乎鲁。\\
\hspace*{2em}孔子年四十二。鲁昭公卒於乾侯。定公立。定公立五年。夏。季平子卒。桓子嗣立。季桓子穿井得土缶。中若羊。问仲尼云﹁得狗﹂。仲尼曰。﹁以丘所闻。羊也。丘闻之。木石之怪夔。罔阆。水之怪龙。罔象。土之怪坟羊。﹂\\
\hspace*{2em}吴伐越。堕会稽。得骨节专车。吴使使问仲尼。﹁骨何者最大。﹂仲尼曰。﹁禹致群神於会稽山。防风氏後至。禹杀而戮之。其节专车。此为大矣。﹂吴客曰。﹁谁为神。﹂仲尼曰。﹁山川之神足以纲纪天下。其守为神。社稷为公侯。皆属於王者。﹂客曰。﹁防风何守。﹂仲尼曰。﹁汪罔氏之君守封。禺之山。为釐姓。在虞。夏。商为汪罔。於周为长翟。今谓之大人。﹂客曰。﹁人长几何。﹂仲尼曰。﹁僬侥氏三尺。短之至也。长者不过十之。数之极也。﹂於是吴客曰。﹁善哉圣人。﹂\\
\hspace*{2em}桓子嬖臣曰仲梁怀。与阳虎有隙。阳虎欲逐怀。公山不狃止之。其秋。怀益骄。阳虎执怀。桓子怒。阳虎因囚桓子。与盟而醳之。阳虎由此益轻季氏。季氏亦僭於公室。陪臣执国政。是以鲁自大夫以下皆僭离於正道。故孔子不仕。退而脩诗书礼乐。弟子弥众。至自远方。莫不受业焉。\\
\hspace*{2em}定公八年。公山不狃不得意於季氏。因阳虎为乱。欲废三桓之適。更立其庶孽阳虎素所善者。遂执季桓子。桓子诈之。得脱。定公九年。阳虎不胜。奔于齐。是时孔子年五十。\\
\hspace*{2em}公山不狃以费畔季氏。使人召孔子。孔子循道弥久。温温无所试。莫能己用。曰。﹁盖周文武起丰镐而王。今费虽小。傥庶几乎。﹂欲往。子路不说。止孔子。孔子曰。﹁夫召我者岂徒哉。如用我。其为东周乎。﹂然亦卒不行。\\
\hspace*{2em}其後定公以孔子为中都宰。一年。四方皆则之。由中都宰为司空。由司空为大司寇。\\
\hspace*{2em}定公十年春。及齐平。夏。齐大夫黎鉏言於景公曰。﹁鲁用孔丘。其势危齐。﹂乃使使告鲁为好会。会於夹谷。鲁定公且以乘车好往。孔子摄相事。曰。﹁臣闻有文事者必有武备。有武事者必有文备。古者诸侯出疆。必具官以从。请具左右司马。﹂定公曰。﹁诺。﹂具左右司马。会齐侯夹谷。为坛位。土阶三等。以会遇之礼相见。揖让而登。献酬之礼毕。齐有司趋而进曰。﹁请奏四方之乐。﹂景公曰。﹁诺。﹂於是旍旄羽袚矛戟剑拨鼓噪而至。孔子趋而进。历阶而登。不尽一等。举袂而言曰。﹁吾两君为好会。夷狄之乐何为於此。请命有司。﹂有司卻之。不去。则左右视晏子与景公。景公心怍。麾而去之。有顷。齐有司趋而进曰。﹁请奏宫中之乐。﹂景公曰。﹁诺。﹂优倡侏儒为戏而前。孔子趋而进。历阶而登。不尽一等。曰。﹁匹夫而营惑诸侯者罪当诛。请命有司。﹂有司加法焉。手足异处。景公惧而动。知义不若。归而大恐。告其群臣曰。﹁鲁以君子之道辅其君。而子独以夷狄之道教寡人。使得罪於鲁君。为之柰何。﹂有司进对曰。﹁君子有过则谢以质。小人有过则谢以文。君若悼之。则谢以质。﹂於是齐侯乃归所侵鲁之郓。汶阳。龟阴之田以谢过。\\
\hspace*{2em}定公十三年夏。孔子言於定公曰。﹁臣无藏甲。大夫毋百雉之城。﹂使仲由为季氏宰。将堕三都。於是叔孙氏先堕郈。季氏将堕费。公山不狃。叔孙辄率费人袭鲁。公与三子入于季氏之宫。登武子之台。费人攻之。弗克。入及公侧。孔子命申句须。乐颀下伐之。费人北。国人追之。败诸姑蔑。二子奔齐。遂堕费。将堕成。公敛处父谓孟孙曰。﹁堕成。齐人必至于北门。且成。孟氏之保鄣。无成是无孟氏也。我将弗堕。﹂十二月。公围成。弗克。\\
\hspace*{2em}定公十四年。孔子年五十六。由大司寇行摄相事。有喜色。门人曰。﹁闻君子祸至不惧。福至不喜。﹂孔子曰。﹁有是言也。不曰﹁乐其以贵下人﹂乎。﹂於是诛鲁大夫乱政者少正卯。与闻国政三月。粥羔豚者弗饰贾。男女行者别於涂。涂不拾遗。四方之客至乎邑者不求有司。皆予之以归。\\
\hspace*{2em}齐人闻而惧。曰。﹁孔子为政必霸。霸则吾地近焉。我之为先并矣。盍致地焉。﹂黎鉏曰。﹁请先尝沮之。沮之而不可则致地。庸迟乎。﹂於是选齐国中女子好者八十人。皆衣文衣而舞康乐。文马三十驷。遗鲁君。陈女乐文马於鲁城南高门外。季桓子微服往观再三。将受。乃语鲁君为周道游。往观终日。怠於政事。子路曰。﹁夫子可以行矣。﹂孔子曰。﹁鲁今且郊。如致膰乎大夫。则吾犹可以止。﹂桓子卒受齐女乐。三日不听政。郊。又不致膰俎於大夫。孔子遂行。宿乎屯。而师己送。曰。﹁夫子则非罪。﹂孔子曰。﹁吾歌可夫。﹂歌曰。﹁彼妇之口。可以出走。彼妇之谒。可以死败。盖优哉游哉。维以卒岁。﹂师己反。桓子曰。﹁孔子亦何言。﹂师己以实告。桓子喟然叹曰。﹁夫子罪我以群婢故也夫。﹂\\
\hspace*{2em}孔子遂適卫。主於子路妻兄颜浊邹家。卫灵公问孔子。﹁居鲁得禄几何。﹂对曰。﹁奉粟六万。﹂卫人亦致粟六万。居顷之。或谮孔子於卫灵公。灵公使公孙余假一出一入。孔子恐获罪焉。居十月。去卫。\\
\hspace*{2em}将適陈。过匡。颜刻为仆。以其策指之曰。﹁昔吾入此。由彼缺也。﹂匡人闻之。以为鲁之阳虎。阳虎尝暴匡人。匡人於是遂止孔子。孔子状类阳虎。拘焉五日。颜渊後。子曰。﹁吾以汝为死矣。﹂颜渊曰。﹁子在。回何敢死。﹂匡人拘孔子益急。弟子惧。孔子曰。﹁文王既没。文不在兹乎。天之将丧斯文也。後死者不得与于斯文也。天之未丧斯文也。匡人其如予何。﹂孔子使从者为甯武子臣於卫。然後得去。\\
\hspace*{2em}去即过蒲。月馀。反乎卫。主蘧伯玉家。灵公夫人有南子者。使人谓孔子曰。﹁四方之君子不辱欲与寡君为兄弟者。必见寡小君。寡小君原见。﹂孔子辞谢。不得已而见之。夫人在絺帷中。孔子入门。北面稽首。夫人自帷中再拜。环珮玉声璆然。孔子曰。﹁吾乡为弗见。见之礼答焉。﹂子路不说。孔子矢之曰。﹁予所不者。天厌之。天厌之。﹂居卫月馀。灵公与夫人同车。宦者雍渠参乘。出。使孔子为次乘。招摇巿过之。孔子曰。﹁吾未见好德如好色者也。﹂於是丑之。去卫。过曹。是岁。鲁定公卒。\\
\hspace*{2em}孔子去曹適宋。与弟子习礼大树下。宋司马桓魋欲杀孔子。拔其树。孔子去。弟子曰。﹁可以速矣。﹂孔子曰。﹁天生德於予。桓魋其如予何。﹂\\
\hspace*{2em}孔子適郑。与弟子相失。孔子独立郭东门。郑人或谓子贡曰。﹁东门有人。其颡似尧。其项类皋陶。其肩类子产。然自要以下不及禹三寸。累累若丧家之狗。﹂子贡以实告孔子。孔子欣然笑曰。﹁形状。末也。而谓似丧家之狗。然哉。然哉。﹂\\
\hspace*{2em}孔子遂至陈。主於司城贞子家。岁馀。吴王夫差伐陈。取三邑而去。赵鞅伐朝歌。楚围蔡。蔡迁于吴。吴败越王勾践会稽。\\
\hspace*{2em}有隼集于陈廷而死。楛矢贯之。石砮。矢长尺有咫。陈湣公使使问仲尼。仲尼曰。﹁隼来远矣。此肃慎之矢也。昔武王克商。通道九夷百蛮。使各以其方贿来贡。使无忘职业。於是肃慎贡楛矢石砮。长尺有咫。先王欲昭其令德。以肃慎矢分大姬。配虞胡公而封诸陈。分同姓以珍玉。展亲。分异姓以远职。使无忘服。故分陈以肃慎矢。﹂试求之故府。果得之。\\
\hspace*{2em}孔子居陈三岁。会晋楚争彊。更伐陈。及吴侵陈。陈常被寇。孔子曰。﹁归与归与。吾党之小子狂简。进取不忘其初。﹂於是孔子去陈。\\
\hspace*{2em}过蒲。会公叔氏以蒲畔。蒲人止孔子。弟子有公良孺者。以私车五乘从孔子。其为人长贤。有勇力。谓曰。﹁吾昔从夫子遇难於匡。今又遇难於此。命也已。吾与夫子再罹难。宁斗而死。﹂斗甚疾。蒲人惧。谓孔子曰。﹁苟毋適卫。吾出子。﹂与之盟。出孔子东门。孔子遂適卫。子贡曰。﹁盟可负邪。﹂孔子曰。﹁要盟也。神不听。﹂\\
\hspace*{2em}卫灵公闻孔子来。喜。郊迎。问曰。﹁蒲可伐乎。﹂对曰。﹁可。﹂灵公曰。﹁吾大夫以为不可。今蒲。卫之所以待晋楚也。以卫伐之。无乃不可乎。﹂孔子曰。﹁其男子有死之志。妇人有保西河之志。吾所伐者不过四五人。﹂灵公曰。﹁善。﹂然不伐蒲。\\
\hspace*{2em}灵公老。怠於政。不用孔子。孔子喟然叹曰。﹁苟有用我者。期月而已。三年有成。﹂孔子行。\\
\hspace*{2em}佛肸为中牟宰。赵简子攻范。中行。伐中牟。佛肸畔。使人召孔子。孔子欲往。子路曰。﹁由闻诸夫子。﹁其身亲为不善者。君子不入也﹂。今佛肸亲以中牟畔。子欲往。如之何。﹂孔子曰。﹁有是言也。不曰坚乎。磨而不磷。不曰白乎。涅而不淄。我岂匏瓜也哉。焉能系而不食。﹂\\
\hspace*{2em}孔子击磬。有荷蒉而过门者。曰。﹁有心哉。击磬乎。硁硁乎。莫己知也夫而已矣。﹂\\
\hspace*{2em}孔子学鼓琴师襄子。十日不进。师襄子曰。﹁可以益矣。﹂孔子曰。﹁丘已习其曲矣。未得其数也。﹂有间。曰。﹁已习其数。可以益矣。﹂孔子曰。﹁丘未得其志也。﹂有间。曰。﹁已习其志。可以益矣。﹂孔子曰。﹁丘未得其为人也。﹂有间。有所穆然深思焉。有所怡然高望而远志焉。曰。﹁丘得其为人。黯然而黑。几然而长。眼如望羊。如王四国。非文王其谁能为此也。﹂师襄子辟席再拜。曰。﹁师盖云文王操也。﹂\\
\hspace*{2em}孔子既不得用於卫。将西见赵简子。至於河而闻窦鸣犊。舜华之死也。临河而叹曰。﹁美哉水。洋洋乎。丘之不济此。命也夫。﹂子贡趋而进曰。﹁敢问何谓也。﹂孔子曰。﹁窦鸣犊。舜华。晋国之贤大夫也。赵简子未得志之时。须此两人而后从政。及其已得志。杀之乃从政。丘闻之也。刳胎杀夭则麒麟不至郊。竭泽涸渔则蛟龙不合阴阳。覆巢毁卵则凤皇不翔。何则。君子讳伤其类也。夫鸟兽之於不义也尚知辟之。而况乎丘哉。﹂乃还息乎陬乡。作为陬操以哀之。而反乎卫。入主蘧伯玉家。\\
\hspace*{2em}他日。灵公问兵陈。孔子曰。﹁俎豆之事则尝闻之。军旅之事未之学也。﹂明日。与孔子语。见蜚雁。仰视之。色不在孔子。孔子遂行。复如陈。\\
\hspace*{2em}夏。卫灵公卒。立孙辄。是为卫出公。六月。赵鞅内太子蒯聩于戚。阳虎使太子絻。八人衰绖。伪自卫迎者。哭而入。遂居焉。冬。蔡迁于州来。是岁鲁哀公三年。而孔子年六十矣。齐助卫围戚。以卫太子蒯聩在故也。\\
\hspace*{2em}夏。鲁桓釐庙燔。南宫敬叔救火。孔子在陈。闻之。曰。﹁灾必於桓釐庙乎。﹂已而果然。\\
\hspace*{2em}秋。季桓子病。辇而见鲁城。喟然叹曰。﹁昔此国几兴矣。以吾获罪於孔子。故不兴也。﹂顾谓其嗣康子曰。﹁我即死。若必相鲁。相鲁。必召仲尼。﹂后数日。桓子卒。康子代立。已葬。欲召仲尼。公之鱼曰。﹁昔吾先君用之不终。终为诸侯笑。今又用之。不能终。是再为诸侯笑。﹂康子曰。﹁则谁召而可。﹂曰。﹁必召冉求。﹂於是使使召冉求。冉求将行。孔子曰。﹁鲁人召求。非小用之。将大用之也。﹂是日。孔子曰。﹁归乎归乎。吾党之小子狂简。斐然成章。吾不知所以裁之。﹂子赣知孔子思归。送冉求。因诫曰﹁即用。以孔子为招﹂云。\\
\hspace*{2em}冉求既去。明年。孔子自陈迁于蔡。蔡昭公将如吴。吴召之也。前昭公欺其臣迁州来。後将往。大夫惧复迁。公孙翩射杀昭公。楚侵蔡。秋。齐景公卒。\\
\hspace*{2em}明年。孔子自蔡如叶。叶公问政。孔子曰。﹁政在来远附迩。﹂他日。叶公问孔子於子路。子路不对。孔子闻之。曰。﹁由。尔何不对曰﹁其为人也。学道不倦。诲人不厌。发愤忘食。乐以忘忧。不知老之将至﹂云尔。﹂\\
\hspace*{2em}去叶。反于蔡。长沮。桀溺耦而耕。孔子以为隐者。使子路问津焉。长沮曰。﹁彼执舆者为谁。﹂子路曰。﹁为孔丘。﹂曰。﹁是鲁孔丘与。﹂曰。﹁然。﹂曰。﹁是知津矣。﹂桀溺谓子路曰。﹁子为谁。﹂曰。﹁为仲由。﹂曰。﹁子。孔丘之徒与。﹂曰。﹁然。﹂桀溺曰。﹁悠悠者天下皆是也。而谁以易之。且与其从辟人之士。岂若从辟世之士哉。﹂櫌而不辍。子路以告孔子。孔子怃然曰。﹁鸟兽不可与同群。天下有道。丘不与易也。﹂\\
\hspace*{2em}他日。子路行。遇荷筱丈人。曰。﹁子见夫子乎。﹂丈人曰。﹁四体不勤。五穀不分。孰为夫子。﹂植其杖而芸。子路以告。孔子曰。﹁隐者也。﹂复往。则亡。\\
\hspace*{2em}孔子迁于蔡三岁。吴伐陈。楚救陈。军于城父。闻孔子在陈蔡之间。楚使人聘孔子。孔子将往拜礼。陈蔡大夫谋曰。﹁孔子贤者。所刺讥皆中诸侯之疾。今者久留陈蔡之间。诸大夫所设行皆非仲尼之意。今楚。大国也。来聘孔子。孔子用於楚。则陈蔡用事大夫危矣。﹂於是乃相与发徒役围孔子於野。不得行。绝粮。从者病。莫能兴。孔子讲诵弦歌不衰。子路愠见曰。﹁君子亦有穷乎。﹂孔子曰。﹁君子固穷。小人穷斯滥矣。﹂\\
\hspace*{2em}子贡色作。孔子曰。﹁赐。尔以予为多学而识之者与。﹂曰。﹁然。非与。﹂孔子曰。﹁非也。予一以贯之。﹂\\
\hspace*{2em}孔子知弟子有愠心。乃召子路而问曰。﹁诗云﹁匪兕匪虎。率彼旷野﹂。吾道非邪。吾何为於此。﹂子路曰。﹁意者吾未仁邪。人之不我信也。意者吾未知邪。人之不我行也。﹂孔子曰。﹁有是乎。由。譬使仁者而必信。安有伯夷。叔齐。使知者而必行。安有王子比干。﹂\\
\hspace*{2em}子路出。子贡入见。孔子曰。﹁赐。诗云﹁匪兕匪虎。率彼旷野﹂。吾道非邪。吾何为於此。﹂子贡曰。﹁夫子之道至大也。故天下莫能容夫子。夫子盖少贬焉。﹂孔子曰。﹁赐。良农能稼而不能为穑。良工能巧而不能为顺。君子能脩其道。纲而纪之。统而理之。而不能为容。今尔不脩尔道而求为容。赐。而志不远矣。﹂\\
\hspace*{2em}子贡出。颜回入见。孔子曰。﹁回。诗云﹁匪兕匪虎。率彼旷野﹂。吾道非邪。吾何为於此。﹂颜回曰。﹁夫子之道至大。故天下莫能容。虽然。夫子推而行之。不容何病。不容然後见君子。夫道之不脩也。是吾丑也。夫道既已大脩而不用。是有国者之丑也。不容何病。不容然後见君子。﹂孔子欣然而笑曰。﹁有是哉颜氏之子。使尔多财。吾为尔宰。﹂\\
\hspace*{2em}於是使子贡至楚。楚昭王兴师迎孔子。然後得免。\\
\hspace*{2em}昭王将以书社地七百里封孔子。楚令尹子西曰。﹁王之使使诸侯有如子贡者乎。﹂曰。﹁无有。﹂﹁王之辅相有如颜回者乎。﹂曰。﹁无有。﹂﹁王之将率有如子路者乎。﹂曰。﹁无有。﹂﹁王之官尹有如宰予者乎。﹂曰。﹁无有。﹂﹁且楚之祖封於周。号为子男五十里。今孔丘述三五之法。明周召之业。王若用之。则楚安得世世堂堂方数千里乎。夫文王在丰。武王在镐。百里之君卒王天下。今孔丘得据土壤。贤弟子为佐。非楚之福也。﹂昭王乃止。其秋。楚昭王卒于城父。\\
\hspace*{2em}楚狂接舆歌而过孔子。曰。﹁凤兮凤兮。何德之衰。往者不可谏兮。来者犹可追也。已而已而。今之从政者殆而。﹂孔子下。欲与之言。趋而去。弗得与之言。\\
\hspace*{2em}於是孔子自楚反乎卫。是岁也。孔子年六十三。而鲁哀公六年也。\\
\hspace*{2em}其明年。吴与鲁会缯。徵百牢。太宰嚭召季康子。康子使子贡往。然后得已。\\
\hspace*{2em}孔子曰。﹁鲁卫之政。兄弟也。﹂是时。卫君辄父不得立。在外。诸侯数以为让。而孔子弟子多仕於卫。卫君欲得孔子为政。子路曰。﹁卫君待子而为政。子将奚先。﹂孔子曰。﹁必也正名乎。﹂子路曰。﹁有是哉。子之迂也。何其正也。﹂孔子曰。﹁野哉由也。夫名不正则言不顺。言不顺则事不成。事不成则礼乐不兴。礼乐不兴则刑罚不中。刑罚不中则民无所错手足矣。夫君子为之必可名。言之必可行。君子於其言。无所苟而已矣。﹂\\
\hspace*{2em}其明年。厓有为季氏将师。与齐战於郎。克之。季康子曰。﹁子之於军旅。学之乎。性之乎。﹂厓有曰。﹁学之於孔子。﹂季康子曰。﹁孔子何如人哉。﹂对曰。﹁用之有名。播之百姓。质诸鬼神而无憾。求之至於此道。虽累千社。夫子不利也。﹂康子曰。﹁我欲召之。可乎。﹂对曰。﹁欲召之。则毋以小人固之。则可矣。﹂而卫孔文子将攻太叔。问策於仲尼。仲尼辞不知。退而命载而行。曰。﹁鸟能择木。木岂能择鸟乎。﹂文子固止。会季康子逐公华。公宾。公林。以币迎孔子。孔子归鲁。\\
\hspace*{2em}孔子之去鲁凡十四岁而反乎鲁。\\
\hspace*{2em}鲁哀公问政。对曰。﹁政在选臣。﹂季康子问政。曰。﹁举直错诸枉。则枉者直。﹂康子患盗。孔子曰。﹁苟子之不欲。虽赏之不窃。﹂然鲁终不能用孔子。孔子亦不求仕。\\
\hspace*{2em}孔子之时。周室微而礼乐废。诗书缺。追迹三代之礼。序书传。上纪唐虞之际。下至秦缪。编次其事。曰。﹁夏礼吾能言之。杞不足徵也。殷礼吾能言之。宋不足徵也。足。则吾能徵之矣。﹂观殷夏所损益。曰。﹁後虽百世可知也。以一文一质。周监二代。郁郁乎文哉。吾从周。﹂故书传。礼记自孔氏。\\
\hspace*{2em}孔子语鲁大师。﹁乐其可知也。始作翕如。纵之纯如。皦如。绎如也。以成。﹂﹁吾自卫反鲁。然後乐正。雅颂各得其所。﹂\\
\hspace*{2em}古者诗三千馀篇。及至孔子。去其重。取可施於礼义。上采契后稷。中述殷周之盛。至幽厉之缺。始於衽席。故曰﹁关雎之乱以为风始。鹿鸣为小雅始。文王为大雅始。清庙为颂始﹂。三百五篇孔子皆弦歌之。以求合韶武雅颂之音。礼乐自此可得而述。以备王道。成六艺。\\
\hspace*{2em}孔子晚而喜易。序彖。系。象。说卦。文言。读易。韦编三绝。曰。﹁假我数年。若是。我於易则彬彬矣。﹂\\
\hspace*{2em}孔子以诗书礼乐教。弟子盖三千焉。身通六艺者七十有二人。如颜浊邹之徒。颇受业者甚众。\\
\hspace*{2em}孔子以四教。文。行。忠。信。绝四。毋意。毋必。毋固。毋我。所慎。齐。战。疾。子罕言利与命与仁。不愤不启。举一隅不以三隅反。则弗复也。\\
\hspace*{2em}其於乡党。恂恂似不能言者。其於宗庙朝廷。辩辩言。唯谨尔。朝。与上大夫言。訚訚如也。与下大夫言。侃侃如也。\\
\hspace*{2em}入公门。鞠躬如也。趋进。翼如也。君召使儐。色勃如也。君命召。不俟驾行矣。\\
\hspace*{2em}鱼馁。肉败。割不正。不食。席不正。不坐。食於有丧者之侧。未尝饱也。\\
\hspace*{2em}是日哭。则不歌。见齐衰。瞽者。虽童子必变。\\
\hspace*{2em}﹁三人行。必得我师。﹂﹁德之不脩。学之不讲。闻义不能徙。不善不能改。是吾忧也。﹂使人歌。善。则使复之。然后和之。\\
\hspace*{2em}子不语。怪。力。乱。神。\\
\hspace*{2em}子贡曰。﹁夫子之文章。可得闻也。夫子言天道与性命。弗可得闻也已。﹂颜渊喟然叹曰。﹁仰之弥高。钻之弥坚。瞻之在前。忽焉在後。夫子循循然善诱人。博我以文。约我以礼。欲罢不能。既竭我才。如有所立。卓尔。虽欲从之。蔑由也已。﹂达巷党人曰。﹁大哉孔子。博学而无所成名。﹂子闻之曰。﹁我何执。执御乎。执射乎。我执御矣。﹂牢曰。﹁子云﹁不试。故艺﹂。﹂\\
\hspace*{2em}鲁哀公十四年春。狩大野。叔孙氏车子鉏商获兽。以为不祥。仲尼视之。曰。﹁麟也。﹂取之。曰。﹁河不出图。雒不出书。吾已矣夫。﹂颜渊死。孔子曰。﹁天丧予。﹂及西狩见麟。曰。﹁吾道穷矣。﹂喟然叹曰。﹁莫知我夫。﹂子贡曰。﹁何为莫知子。﹂子曰。﹁不怨天。不尤人。下学而上达。知我者其天乎。﹂\\
\hspace*{2em}﹁不降其志。不辱其身。伯夷。叔齐乎。﹂谓﹁柳下惠。少连降志辱身矣﹂。谓﹁虞仲。夷逸隐居放言。行中清。废中权﹂。﹁我则异於是。无可无不可。﹂\\
\hspace*{2em}子曰。﹁弗乎弗乎。君子病没世而名不称焉。吾道不行矣。吾何以自见於後世哉。﹂乃因史记作春秋。上至隐公。下讫哀公十四年。十二公。据鲁。亲周。故殷。运之三代。约其文辞而指博。故吴楚之君自称王。而春秋贬之曰﹁子﹂。践土之会实召周天子。而春秋讳之曰﹁天王狩於河阳﹂。推此类以绳当世。贬损之义。後有王者举而开之。春秋之义行。则天下乱臣贼子惧焉。\\
\hspace*{2em}孔子在位听讼。文辞有可与人共者。弗独有也。至於为春秋。笔则笔。削则削。子夏之徒不能赞一辞。弟子受春秋。孔子曰。﹁後世知丘者以春秋。而罪丘者亦以春秋。﹂\\
\hspace*{2em}明岁。子路死於卫。孔子病。子贡请见。孔子方负杖逍遥於门。曰。﹁赐。汝来何其晚也。﹂孔子因叹。歌曰。﹁太山坏乎。梁柱摧乎。哲人萎乎。﹂因以涕下。谓子贡曰。﹁天下无道久矣。莫能宗予。夏人殡於东阶。周人於西阶。殷人两柱间。昨暮予梦坐奠两柱之间。予始殷人也。﹂後七日卒。\\
\hspace*{2em}孔子年七十三。以鲁哀公十六年四月己丑卒。\\
\hspace*{2em}哀公诔之曰。﹁旻天不吊。不玦遗一老。俾屏余一人以在位。茕茕余在疚。呜呼哀哉。尼父。毋自律。﹂子贡曰。﹁君其不没於鲁乎。夫子之言曰。﹁礼失则昏。名失则愆。失志为昏。失所为愆。﹂生不能用。死而诔之。非礼也。称﹁余一人﹂。非名也。﹂\\
\hspace*{2em}孔子葬鲁城北泗上。弟子皆服三年。三年心丧毕。相诀而去。则哭。各复尽哀。或复留。唯子赣庐於冢上。凡六年。然後去。弟子及鲁人往从冢而家者百有馀室。因命曰孔里。鲁世世相传以岁时奉祠孔子冢。而诸儒亦讲礼乡饮大射於孔子冢。孔子冢大一顷。故所居堂弟子内。後世因庙藏孔子衣冠琴车书。至于汉二百馀年不绝。高皇帝过鲁。以太牢祠焉。诸侯卿相至。常先谒然後从政。\\
\hspace*{2em}孔子生鲤。字伯鱼。伯鱼年五十。先孔子死。\\
\hspace*{2em}伯鱼生伋。字子思。年六十二。尝困於宋。子思作中庸。\\
\hspace*{2em}子思生白。字子上。年四十七。子上生求。字子家。年四十五。子家生箕。字子京。年四十六。子京生穿。字子高。年五十一。子高生子慎。年五十七。尝为魏相。\\
\hspace*{2em}子慎生鲋。年五十七。为陈王涉博士。死於陈下。\\
\hspace*{2em}鲋弟子襄。年五十七。尝为孝惠皇帝博士。迁为长沙太守。长九尺六寸。\\
\hspace*{2em}子襄生忠。年五十七。忠生武。武生延年及安国。安国为今皇帝博士。至临淮太守。蚤卒。安国生卬。卬生驩。\\
\hspace*{2em}太史公曰。诗有之。﹁高山仰止。景行行止。﹂虽不能至。然心乡往之。余读孔氏书。想见其为人。適鲁。观仲尼庙堂车服礼器。诸生以时习礼其家。余祗回留之不能去云。天下君王至于贤人众矣。当时则荣。没则已焉。孔子布衣。传十馀世。学者宗之。自天子王侯。中国言六艺者折中於夫子。可谓至圣矣。\\

% ===== 索隐述赞（司马贞《史记索隐》）· 统一样式 =====
% 题头「索隐述赞曰」五字并入反白框；正文每四句一列，定宽盒 \hfil 居列中；
% \fontsize{24}{41} 列距对齐全书界行（121+41k）；\clearpage 使述赞独占一页。
\clearpage
\shuzanhead{索隐述赞曰}\linebreak
{\shuzanbody\fontsize{24}{41}\selectfont
\hbox to \hsize{\hfil 孔子之胄。出于商国。弗父能让。正考铭勒。\hfil}\linebreak
\hbox to \hsize{\hfil 防叔来奔。邹人掎足。尼丘诞圣。阙里生德。\hfil}\linebreak
\hbox to \hsize{\hfil 七十升堂。四方取则。卯诛两观。摄相夹谷。\hfil}\linebreak
\hbox to \hsize{\hfil 歌凤遽衰。泣麟何促。九流仰镜。万古钦躅。\hfil}%
\par}
\end{document}
